%--- Preamble ---------------------------------------------------------%
% Load LaTeX packages
\documentclass[11pt, xcolor=dvipsnames]{beamer}
\usepackage[absolute, overlay]{textpos}                           % supports floating text in any location
\usepackage[utf8]{inputenc}
\usepackage{csquotes,xpatch}% recommended
%\usepackage[english]{babel}
%\usepackage[american]{babel}
\usepackage[brazil]{babel}
\usepackage{hyperref}
\usepackage{graphicx}
\usepackage{tabularx}
\usepackage{booktabs}
\usepackage{bookmark}
\usepackage{pgfplots}
\usepackage{tikz}

% Customize theme attributes
\useoutertheme{infolines}					                     % adds 3 box footer
\usetheme[height=7mm]{Rochester} 				                 % choose theme
\setbeamertemplate{blocks}[rounded][shadow=true] 		         % rounded theorem box with shadow
\setbeamertemplate{caption}[numbered]                             % enable counting of tables/figures

% Define template colors
\definecolor{QPblue}{RGB}{0,25,100}                               % define QP blue using RGB code
\definecolor{QPgreen}{RGB}{0,153,110}                             % define QP green using RGB code
\setbeamercolor{title}{fg=white, bg=QPblue}                       % define title page box color
\setbeamercolor{frametitle}{fg=white, bg=QPblue}                  % define frame title color
\setbeamercolor{normal text}{fg=black}                            % define standard font color
\setbeamercolor{author in head/foot}{fg=QPblue, bg=QPblue!75}    % define infoline 1st box color
\setbeamercolor{title in head/foot}{fg=QPblue, bg=QPblue!60}     % define info line 2nd box color
\setbeamercolor{date in head/foot}{fg=QPblue, bg=QPblue!30}	  % define infoline 3rd box color
\setbeamercolor{block title}{fg=QPblue!50!QPgreen, bg=QPblue!30} % define theorem box title color
\setbeamercolor{block body}{fg=QPblue!50!QPgreen, bg=gray!10}	 % define theorem box body color
\setbeamercolor{local structure}{fg=QPblue!75}		           % define bullet and enumerate list colors

% Define global environments
\newenvironment{reference}[2]{                                    % define environment for footnotes
  \begin{textblock*}{\textwidth}(#1, #2)
      \tiny\it\bgroup\color{red!70!QPblue}}{\egroup\end{textblock*}}

% Define title page logo and project metadata
\titlegraphic{\includegraphics[width=3cm]{../images/uninove.png}\hspace*{0cm}~
}

\title{Estatística Bayesiana}
\subtitle{Distribuições Probabilísticas}
\author{Jose Storopoli}
\institute[PPGA/CIS]{
   \textcolor{QPblue!75}{Universidade Nove de Julho \\
   UNINOVE \\
   São Paulo \\
   Brasil \\ [1ex]
   \texttt{josees@uni9.pro.br}} \\ [1ex]
   Creative Commons Attribution-ShareAlike 4.0
}
\date{Maio 2022}

\begin{document}

%--- Distributions ----------------------------------------------------%

\tikzset{
    declare function={
      gauss(\m,\s)=1/(\s*sqrt(2*pi))*exp(-((x-\m)^2)/(2*\s^2));%
    }
}
%--- Title Page -------------------------------------------------------%

\begin{frame}[plain]
  \titlepage
  \begin{reference}{4mm}{90mm}
    XXX
  \end{reference}
\end{frame}

%--- Slide 1 ----------------------------------------------------------%

\begin{frame}{Normal}
    \centering
    \begin{tikzpicture}[line join=round, line cap=round]
        \begin{axis}[every axis plot post/.append style={
            mark=none,domain=-4:4,samples=200,smooth}, % All plots: from -2:2, 50 samples, smooth, no marks
            axis x line*=bottom, % no box around the plot, only x and y axis
            axis y line*=left, % the * suppresses the arrow tips
            enlargelimits=upper,
            %restrict y to domain=0:10
            ] % extend the axes a bit to the right and top

            \addplot [red] {gauss(0, 1)};
            \addlegendentry{$\mu=0, \sigma=1$}
            \addplot [blue] {gauss(0, 2)};
            \addlegendentry{$\mu=0, \sigma=2$}
            %\legend{$\mu=0, \sigma=1$, $\mu=0, \sigma=2$};
        \end{axis}
    \end{tikzpicture}
\end{frame}

%--- End -------------------------------------------------------------%
\end{document}
